\section{Cơ hội khởi nghiệp}

\begin{frame}{Esports Việt Nam bước vào giai đoạn chuyên nghiệp hóa}{Lễ ký kết chiến lược 2026--2030 là tín hiệu cho một hệ sinh thái đang mở rộng}
  \begin{columns}[c,onlytextwidth]
    \column{.48\textwidth}
      \includegraphics[width=\linewidth,height=.56\textheight,keepaspectratio]{clavis-esport-inline-1.jpg}\par
      \vspace{.3em}
      {\centering\scriptsize\color{cookieMuted}
      Lễ ký kết chiến lược phát triển Esports quốc gia 2026--2030, ngày 14/5/2026.\par}
    \column{.43\textwidth}
      \cookiekicker{TÍN HIỆU THỊ TRƯỜNG}\par\vspace{.45em}
      \begin{itemize}
        \item Chiến lược được ký kết ngày 14/5/2026.
        \item Hệ sinh thái được định hướng phát triển qua công nghệ, giải đấu, doanh nghiệp và cộng đồng.
      \end{itemize}
  \end{columns}
  \note{[Sources]
    Nội dung và ảnh: Ngô Võ Quốc Vương, ``Esport là gì? Vì sao Esport đang trở thành ngành kinh tế sáng tạo tại Việt Nam?'', Clavis Academy, 15/05/2026. URL: https://clavis.edu.vn/tin-tuc/the-thao/esport-la-gi-vi-sao-esport-dang-tro-thanh-nganh-kinh-te-sang-tao-tai-viet-nam
  }
\end{frame}

\begin{frame}{Từ xu hướng đến cơ hội cho Mutux}{Cộng đồng Esports tạo thêm điểm chạm để kiểm chứng nhu cầu thuê thử gear}
  \vspace{1em}
  \centering
  {\Large\textcolor{mutuxBlue}{\textbf{Cộng đồng \& giải đấu}}}
  \hspace{.8em}{\color{mutuxOrange}\Large$\Longrightarrow$}\hspace{.8em}
  {\Large\textcolor{mutuxBlue}{\textbf{Trải nghiệm gear}}}
  \hspace{.8em}{\color{mutuxOrange}\Large$\Longrightarrow$}\hspace{.8em}
  {\Large\textcolor{mutuxBlue}{\textbf{Đơn thuê thực}}}\par
  \vfill
  \centering\textcolor{mutuxOrange}{\textbf{Mục tiêu: xác nhận nhu cầu thật trước khi mở rộng.}}
\end{frame}

\begin{frame}{Khảo sát 1/3: Tần suất chơi game}{20/23 người chơi ít nhất 1--2 lần mỗi tuần}
  \centering
  \includegraphics[width=.78\textwidth,height=.58\textheight,keepaspectratio]{frequency_game.png}\par
  \vfill
  \centering\textcolor{mutuxOrange}{\textbf{86,9\% có tần suất chơi hàng tuần: gaming là một hành vi lặp lại.}}
\end{frame}

\begin{frame}{Khảo sát 2/3: Gear đã từng mua}{Chuột, bàn phím và tai nghe là ba nhóm phổ biến nhất}
  \centering
  \includegraphics[width=.82\textwidth,height=.58\textheight,keepaspectratio]{gaming_gear_bought.png}\par
  \vfill
  \centering\textcolor{mutuxOrange}{\textbf{19/23 người đã từng mua gear: thị trường không cần giáo dục từ đầu.}}
\end{frame}

\begin{frame}{Khảo sát 3/3: Mức độ quan tâm thuê thử}{18/23 người hứng thú hoặc sẵn sàng cân nhắc}
  \centering
  \includegraphics[width=.78\textwidth,height=.58\textheight,keepaspectratio]{Maybe_rent.png}\par
  \vfill
  \centering\textcolor{mutuxOrange}{\textbf{78,2\% mở với đề xuất thuê để trải nghiệm trước khi mua.}}
\end{frame}

\begin{frame}{Cách nhóm thí điểm cơ hội này}{Bắt đầu quanh ĐHQG-HCM với phạm vi đủ nhỏ để đo được}
  \vspace{.8em}
  \begin{columns}[T,onlytextwidth]
    \column{.30\textwidth}
      \cookiekicker{Đối tác}\par\vspace{.45em}
      3--5 shop neo và câu lạc bộ sinh viên quanh ĐHQG-HCM.
    \column{.30\textwidth}
      \cookiekicker{Kênh thử nghiệm}\par\vspace{.45em}
      Giải đấu sinh viên, cộng đồng gaming và social content.
    \column{.30\textwidth}
      \cookiekicker{Điều cần đo}\par\vspace{.45em}
      Tỷ lệ tạo đơn, hoàn tất, tranh chấp và tái thuê.
  \end{columns}
\end{frame}
